\documentclass[a4paper]{article}
\usepackage{geometry}
%\geometry{hmargin={2cm,2cm}}
\usepackage{graphicx}
\usepackage{xcolor}
\usepackage{hyperref}
\definecolor{logo}{RGB}{20,99,163}
\usepackage[T1]{fontenc}
\usepackage[utf8]{inputenc}
\usepackage[english,italian]{babel}
\begin{document}
\begin{titlepage}
    \begin{center}
        \includegraphics{../Images/logo}\\
        \vspace{20px}
        \textcolor{logo}{\hrulefill}\\
        \vspace{20px}
        \textbf{\huge\textcolor{logo}{Studio di fattibilità}}\\
        \vspace{10px}
        \textcolor{logo}{\hrulefill}\\
        \vspace{20px}
        \href{mailto:techsweave@gmail.com}{techsweave@gmail.com}\\
        \vspace{40px}
        \textbf{Informazioni sul documento}\\
        \vspace{20px}
        \begin{tabular}{p{100px} | p{100px}}
            \textbf{Versione} & 0.1\\
            \textbf{Responsabile} & \\
            \textbf{Redattore} & Nicolò Giaccone\\
            \textbf{Verificatore} & \\
            \textbf{Destinatari} & Prof. Tullio Vardanega \newline Prof. Riccardo Cardin, \newline TechSWEave\\
            \textbf{Stato} & Da approvare\\
            \textbf{Uso} & Interno\\
        \end{tabular}
            \vspace{1ex}

    \end{center}
    \end{titlepage}
    \begin{center}
       \textbf{Diario delle modifiche}
        \begin{itemize}
            \item \textit{Prima stesura del documento}, Nicolò Giaccone,(ruolo da definire), v0.1;
        \end{itemize}
    \end{center}
    \newpage
\tableofcontents
\newpage
\section{Capitolato scelto: C2 -  	EmporioLambda}
\section{Capitolato C1 - BlockCOVID}
\section{Capitolato C3 - Gathering Detection Platform}
\section{Capitolato C4 - HD Viz}
\section{Capitolato C5 - PORTACS}
\subsection{Informaizoni genarli}
\begin{itemize}
    \item \textbf{Nome:} \textit{PORTACS};
    \item \textbf{Proponente:} \textit{Sanmarco Informatica};
    \item \textbf{Commitenti:} \textit{Prof. Tullio Vardanega,Prof. Riccardo Cardin};
\end{itemize}
\subsection{Descrizione del capitolato}
Il capitolato richiede di un software in grado di gestire un ente centrale che riceve dati 
di posizioni e velocita da varie unità e ne regola gli spostementi. L'applicazione avrà una \textit{UI} che inicherà 
la posizione, la velocita di ogniu unità e la direzione che questa dovrà prendere.
\subsection{Finalità del progetto}
La finalità del progetto è quella di realizzare una applicazione in grado di gestire, lo spostamento di un unità all'interno di un area controllata,
attraverso la consocenza di alcuni dati e con l'obbiettivo di minimizzare costi e tempi di viaggio e di evitare eventuali collisioni
\subsection{Tecnologie interessate}
\begin{itemize}
    \item \textbf{Docker:} progetto open-source che automatizza il deploymentGdi applicazioni all’in-terno di container.
\end{itemize}
\subsection{Aspetti positivi}
\begin{itemize}
    \item L'obbiettivo del progetto è in linea con le problematiche tecnologiche attuali;
    \item Ha una grande valenza formativa per i componenti del gruppo
\end{itemize}
\subsection{Criticità e fattori di rischio}
\begin{itemize}
    \item Il capitolato non presentatava più posti disponibili;
\end{itemize}
\subsection{Conclusioni}
Il capitolato per quanto di forte interesse e utilità non presenta posti disponibili e non è stato quindi possibile sceglierlo.
\section{Capitolato C6 - Realtime Gaming Platform}
\section{Capitolato C7 - soluzioni di sincronizzazione desktop}
\end{document}
